\documentclass[11pt,a4paper]{article}

\usepackage[T2A]{fontenc}
\usepackage[utf8]{inputenc}
\usepackage[english,russian]{babel}
\usepackage{lmodern}
\usepackage{amsmath,amssymb}
\usepackage{xcolor}
\usepackage[margin=2.2cm]{geometry}
\usepackage{parskip}
\usepackage{titlesec}

\definecolor{primary}{HTML}{1B5E7A}
\definecolor{accent}{HTML}{E07A1F}
\definecolor{lighttext}{HTML}{4A4A4A}

\titleformat{\section}{\Large\bfseries\color{primary}}{}{0pt}{}
\titleformat{\subsection}{\large\color{accent}}{}{0pt}{}
\titlespacing*{\section}{0pt}{*2.4}{*1}
\titlespacing*{\subsection}{0pt}{*1.5}{*0.5}

\newcommand{\slide}[3]{%
  \section*{Слайд #1 — #2\hfill\normalsize\color{lighttext}\textit{(\textasciitilde #3)}}%
}

\setlength{\parindent}{0pt}
\setlength{\parskip}{0.6em}

\title{\vspace{-1cm}\color{primary}\bfseries Текст выступления\\
{\large NFG-SS: распределённая невыпуклая оптимизация без полного участия клиентов}}
\author{}
\date{}

\begin{document}
\maketitle
\thispagestyle{empty}

\vspace{-0.8cm}
{\color{accent}\rule{\linewidth}{0.8mm}}

\vspace{0.4cm}
\textbf{Общий хронометраж:} 10 минут (~9:30 чистой речи + буфер на переключение слайдов).
Темп — спокойный, около 150 слов/минуту. Цифры в скобках — ориентир.

\vspace{0.6cm}

% ============================================================
\slide{1}{Титульный слайд}{20 сек}

Добрый день, уважаемые члены комиссии. Меня зовут Зайнуллин Амир, группа Б05-206д. Тема моей преддипломной работы — локальные методы распределённой оптимизации без подсчёта полного градиента. Метод, который я предлагаю, мы назвали \texttt{NFG-SS}. Научный руководитель — Александр Николаевич Безносиков.

\emph{Кратко о плане доклада:} сначала я обозначу контекст и поставлю задачу, потом покажу, какие два направления существуют в литературе и какой между ними gap. После этого расскажу про сам алгоритм и теоретический результат. В конце — эксперимент на CIFAR-10 и выводы.

% ============================================================
\slide{2}{Цель и задачи}{30 сек}

Цель работы — построить распределённый алгоритм для невыпуклой оптимизации, который одновременно использует similarity между данными клиентов для ускорения и при этом не требует вычисления полного градиента ни в одной точке. Дополнительно — линейная по размерности память на сервере, не растущая с числом клиентов.

Задачи делятся на анализ существующих методов в двух направлениях, разработку собственного алгоритма, доказательство его сходимости в невыпуклом случае и экспериментальное сравнение с baseline'ами.

% ============================================================
\slide{3}{Актуальность}{50 сек}

Зачем нам distributed-сетап. Современные модели и датасеты растут — обучить такое на одной машине нельзя, и распределение по узлам становится обязательным. Но за это надо платить связью: каждый раунд это $d$-мерные векторы между сервером и каждым клиентом. Метрика прогресса в этой области не число итераций, а общее число transmitted vectors.

Что нам даёт similarity. Если данные клиентов из одного распределения, локальные функции близки в моменте — норма разности гессианов ограничена параметром $\delta$. И в реальных задачах $\delta$ намного меньше параметра гладкости $L$.

Проблема в литературе. Есть две линии работ. Similarity-методы — используют $\delta$ для ускорения, но в начале каждой эпохи дёргают всех клиентов за полным градиентом. И no-full-gradient методы — они избегают этого, но не учитывают $\delta$. Их сложность зависит от $L$, и весь выигрыш similarity теряется. Эти линии до сих пор развивались отдельно.

% ============================================================
\slide{4}{Постановка задачи}{45 сек}

Формальная задача — стандартная распределённая finite-sum: минимизируем среднее по узлам, где каждая локальная функция $f_i$ — empirical mean по своему датасету. Один узел — сервер, хранит функцию $f_1$; остальные $n-1$ — клиенты.

Теперь зачем нам нужно \emph{прокс-разложение}. У глобальной функции $f$ нет полезной глобальной гладкости — для нейросетей $L$ либо неизвестен, либо чудовищно велик. И выпуклости тоже нет. То есть мы не можем напрямую делать обычный градиентный спуск и надеяться на классические descent-неравенства.

Решение: раскладываем $f$ на серверную часть $f_1$ и клиентскую разность $(f - f_1)$. Серверная часть полностью доступна на сервере, поэтому её мы обрабатываем через прокс-оператор \emph{точно} — это бесплатно. А клиентскую часть нужно \emph{оценить} стохастически, без полного обхода всех клиентов. Это и есть основная работа нашего алгоритма.

% ============================================================
\slide{5}{Second-order similarity}{40 сек}

Формальное определение: норма разности гессианов локальной функции и среднего ограничена $\delta$ во всех точках.

Это допущение строго слабее, чем $L$-гладкость: при $L$-smoothness $\delta$ автоматически не превосходит $2L$. Откуда вообще берётся $\delta$? Из теории обучения: если локальные датасеты — i.i.d.~семплы из общего распределения размера $m$, то $\delta$ скейлится как $L$ делённое на корень из $m$. В современном ML $m$ велико — отсюда $\delta \ll L$.

Почему это критично. Сложность similarity-методов масштабируется как $\delta$, а не $L$ — это выигрыш фактора $L/\delta$ по communication. Грубо говоря, similarity — это бесплатная структура, и грамотный алгоритм её обязан использовать.

% ============================================================
\slide{6}{Связанная работа: similarity-методы}{40 сек}

Первая линия работ. Начало — DANE Шамира 2014 года, Newton-type для quadratic objectives. Дальнейшие улучшения — AIDE, DISCO, GIANT. Параллельно — SONATA и Accelerated Extragradient Ковалёва 2022, которые достигают оптимального communication-bound $\sqrt{\delta/\mu}$ под логарифмом.

Стохастическая ветвь: SVRP и Catalyzed SVRP Халеда и Цзина — variance reduction с прокс-операторами. И SVRS, AccSVRS Lin et al. 2023 с оптимальной сложностью $n + n^{3/4}\sqrt{\delta/\mu}$ под логарифмом.

Общий минус всей этой группы — в начале каждой эпохи каждый клиент шлёт серверу свой $\nabla f_i$, это $\mathcal{O}(n)$ векторов на каждый refresh. Плюс везде нужен strong convexity, что для нейросетей нереалистично.

% ============================================================
\slide{7}{Связанная работа: no-full-gradient методы}{25 сек}

Вторая линия — методы без full grad. SAG и SAGA — таблица всех градиентов, дорогая память. SARAH и SPIDER — recursive estimator, но периодически дёргают полный градиент. STORM и PAGE — momentum-based или probabilistic switching. И полностью без full grad: ZeroSARAH 2021, Shuffled SARAH Безносикова--Такача 2024, и Medyakov et al. 2025, где anchor заменён на running-mean по random reshuffle.

Общий минус — все методы centralized и не учитывают similarity. Сложность зависит от $L$.

% ============================================================
\slide{8}{Пересечение и gap}{40 сек}

Что находится на пересечении двух линий и почему этого недостаточно. Beznosikov и Takáč 2024 — Shuffled SARAH под similarity, но strongly convex и centralized, не distributed. SILVER, Oko et al.~2024 — non-convex distributed без full grad, но платит $\mathcal{O}(nd)$ памяти на сервере, по вектору на клиента. SPAM, Karagulyan 2024 — similarity-aware без full grad, но в stochastic regime, не использует finite-sum и не даёт оптимальной $1/\varepsilon^2$ сходимости в нашем смысле.

Незакрытая ниша — метод, удовлетворяющий одновременно шести требованиям: non-convex, distributed, finite-sum, similarity-aware, без full gradient, и $\mathcal{O}(d)$ память. NFG-SS — первый такой.

% ============================================================
\slide{9}{Основная идея NFG-SS}{60 сек}

Это центральный слайд алгоритмической части. Мы поддерживаем \emph{два} рекурсивных оценщика. Первый — \emph{SARAH-телескоп} $v_t$: классическая low-variance конструкция через разности градиентов на соседних итерациях. Второй — \emph{running mean} $\tilde v_t$ по \emph{уже посчитанным} в эпохе градиентам, со специальными коэффициентами $(t-1)/t$ и $1/(tb)$.

Ключевое наблюдение: после полного обхода всех клиентов в эпохе running-mean — это усреднение градиентов клиентов в разных точках траектории. Если бы все эти точки совпадали с финальной, это был бы \emph{точно} полный градиент в этой точке. На практике точки разные, но bias ограничен суммой квадратов сдвигов итератов внутри эпохи — это то, что мы можем контролировать через step size.

Главный trick: на следующей эпохе берём $v_0^{(s+1)}$ равным финальному running-mean из предыдущей эпохи. Анкер бесплатный — никакого full gradient не считается ни в одной точке.

% ============================================================
\slide{10}{Алгоритм NFG-SS — псевдокод}{40 сек}

Псевдокод формализует идею. Ключевой момент — шаг 2: $v_0^{(s)}$ это \emph{carry-over running-mean} из эпохи $s-1$, а не свежий полный градиент. Внутри эпохи на каждой итерации параллельно обновляем оба оценщика — running-mean (a) и SARAH-телескоп (b) — и делаем прокс-шаг (c). В конце эпохи финальное значение running-mean становится анкером для следующей эпохи (шаг 5). Полный градиент не считается нигде.

% ============================================================
\slide{11}{Теорема о сходимости}{45 сек}

Теорема даёт два слагаемых в oracle complexity. Первое — $n \delta \Delta_1 / \varepsilon^2$ — это слагаемое типа no-full-grad SARAH из Medyakov et al., но с $\delta$ вместо $L$. Именно здесь сохраняется выигрыш similarity. Второе слагаемое — $n^2 \delta^2 \Delta_2 / (b \varepsilon^2)$ — плата за то, что инициализация эпохи строится по итератам предыдущей.

Особо подчеркну три вещи. Первое: никакого предположения о выпуклости. Второе: никакого вычисления $\nabla f$ ни в одной точке. Третье: сложность зависит от similarity-параметра $\delta$, а не от $L$. Так как в реальных задачах $\delta \ll L$, это означает кратное ускорение по сравнению с любым существующим no-full-grad методом.

% ============================================================
\slide{12}{Сравнительная таблица}{30 сек}

Сводная таблица позиционирует NFG-SS среди конкурентов. SVRS и AccSVRS требуют full grad и strong convexity. Shuffled SARAH без full grad, но всё ещё strongly convex и зависит от $L$ в первом слагаемом. SILVER non-convex без full grad, но платит $\mathcal{O}(nd)$ памятью на сервере. SPAM в stochastic regime с $1/\varepsilon^3$.

NFG-SS — единственный метод с полным набором свойств одновременно.

% ============================================================
\slide{13}{Эксперимент на CIFAR-10 / ResNet-18}{55 сек}

Эмпирическая валидация. CIFAR-10, модель ResNet-18, 11 узлов: один сервер и 10 клиентов с uniform random partition. Сравнение с двумя baseline'ами: SVRS — в точной формулировке Lin et al. 2023, с честным full-gradient anchor в начале каждой эпохи. И FedAvg — в communication-matched режиме, то есть бюджет transmitted vectors одинаков.

На графике — test accuracy в зависимости от числа transmitted vectors. Кривая — медиана по трём random seeds, лента — min-max envelope. NFG-SS обгоняет SVRS примерно после тысячи векторов и удерживает преимущество до конца. Финальная медиана: NFG-SS — около 0.79, SVRS — 0.74, FedAvg — 0.71 при одинаковом бюджете 2000 transmitted vectors.

Важно: ленты NFG-SS и SVRS на финальном плато не пересекаются — разница устойчива по сидам.

% ============================================================
\slide{14}{Заключение}{30 сек}

Подытожу. Предложен метод NFG-SS — distributed variance-reduced алгоритм для невыпуклой оптимизации под second-order similarity. Главный результат: это первый distributed-метод, использующий similarity \emph{без полного участия клиентов} ни в одной точке. Сложность зависит от $\delta$, не от $L$; невыпуклый случай; память на сервере $\mathcal{O}(d)$. Эмпирически dominates SVRS и communication-matched FedAvg.

Направления развития: acceleration через Catalyst-обёртку, снятие предположения о полном покрытии клиентов в эпохе, и расширение на decentralized топологии.

% ============================================================
\slide{15}{Спасибо за внимание}{5 сек}

Спасибо за внимание. Готов ответить на ваши вопросы.

\vspace{1.2cm}
{\color{accent}\rule{\linewidth}{0.4mm}}

\vspace{0.4cm}
\textit{Подсказки для подготовки:}

\begin{itemize}
\item Самые «дорогие» по смыслу слайды — 3 (актуальность), 4 (постановка с прокс-разложением), 8 (gap), 9 (идея), 11 (теорема), 13 (эксперимент). На них стоит не торопиться.
\item Слайды 6, 7 (related work) — пройти средним темпом, не вязнуть в названиях методов; логика группировки видна на слайде сама.
\item На слайде 10 (псевдокод) не зачитывай построчно. Покажи на шаги 2 и 5 — где carry-over и где никакого full grad.
\item На слайде 13 после описания постановки сразу поверни голову к графику и пробеги взглядом снизу вверх: NFG-SS, SVRS, FedAvg.
\item Возможные вопросы: «почему именно $\delta$, а не $L$?» (similarity — структурное предположение, $\delta \ll L$ при i.i.d.~данных). «Зачем второй оценщик $\tilde v$?» (анкер для следующей эпохи бесплатно из уже посчитанных градиентов). «Что если $\delta \approx L$?» (тогда наш метод не лучше обычного no-full-grad — выигрыш именно в режиме $\delta \ll L$).
\end{itemize}

\end{document}
